% Figure for Oscillations, Waves and Optics §A7.1 (single slit of width D at the angle of the first minimum: the ray from the bottom edge travels D sin(theta) = lambda farther than the ray from the top edge, so each ray from the lower half travels lambda/2 farther than its partner a distance D/2 above it, and the pairs cancel).
\begin{tikzpicture}[font=\small, >={Stealth[length=2.2mm]}]
  \def\th{25}
  \fill[black!30] (-0.12,-2.6) rectangle (0.12,-1.2);
  \fill[black!30] (-0.12,1.2) rectangle (0.12,2.6);
  \coordinate (T) at (0.12,1.2);
  \coordinate (M) at (0.12,0);
  \coordinate (B) at (0.12,-1.2);
  \coordinate (T2) at (0.12,0.45);
  \coordinate (B2) at (0.12,-0.75);
  \draw[->, black!55, thin] (-2.5,0.6) -- (-1.3,0.6);
  \draw[->, black!55, thin] (-2.5,-0.6) -- (-1.3,-0.6);
  \node[font=\scriptsize, black!60, anchor=south] at (-1.9,0.65) {plane wave};
  \foreach \p/\L in {T/4.0, T2/4.35, M/4.6, B2/4.95, B/5.2} {\draw[figB, thick, ->] (\p) -- ++(\th:\L);}
  % path differences: foot of perpendicular from T onto the ray from B, and from T onto the ray from M
  \coordinate (FB) at ($(B)+(\th:{2.4*sin(\th)})$);
  \coordinate (FM) at ($(M)+(\th:{1.2*sin(\th)})$);
  \draw[black!55, dashed, thin] (T) -- (FB);
  \draw[figR, line width=1.6pt] (B) -- (FB);
  \draw[figO, line width=1.6pt] (M) -- (FM);
  \node[figR, font=\scriptsize, anchor=north west] at ($(B)+(0.35,-0.15)$) {$D\sin\theta = \lambda$};
  \node[figO, font=\scriptsize] at (0.44,-0.22) {$\tfrac{\lambda}{2}$};
  \draw[<->, black!60, thin] (-0.5,-1.2) -- (-0.5,1.2) node[midway, left, font=\scriptsize] {$D$};
  \draw[black!55, thin] (B) -- ++(1.7,0);
  \draw[black!55, thin] ($(B)+(1.4,0)$) arc (0:\th:1.4);
  \node[font=\scriptsize] at ($(B)+(10:1.65)$) {$\theta$};
  \node[font=\scriptsize, figB, anchor=west] at ($(B)+(\th:5.2)$) {to a distant screen};
\end{tikzpicture}
